\songtitle{C'est dommage}

\setin{2005-09-01}
\tag*{2026}\tag{chanson-francaise}\tag{trip-hop}\tag{french-lyrics}\tag{ricochets}\tag{sophie-durand}\tag{fabian-muller}
% -> ricochets.tex
\begin{notes}
Lyrics by Claude based on \emph{Ricochets} by Carlos Thompson -- chapter 14 (\emph{Agent Perez}): in Sophie Durand's voice, a night with beers on the balcony of Fabián Müller's Bal Harbour apartment, reporting his words: born in the mountains, fascinated by the sea, he could get used to this life, and it is a pity it is not with her. Retrieved from rataflechera's Suno page (V6, 24 November 2026).
\end{notes}
\begin{style}
French chanson over downtempo; soft, breathy female mezzo-soprano singing in French, with a brief powerful full-voice belt in the chorus, tender and resigned with a little smile, first-person narration reported in indirect speech; muted accordion lead, brushed drums, upright bass, felt piano, distant surf and a jet passing high as texture; 76 BPM, slow trip-hop pulse, F major turning to D minor in the chorus; close-mic, unprocessed vocals with natural room reflections, intimate and late-night, restrained dynamics, bass and piano bridge, accordion-only final refrain.
\end{style}
\begin{lyrics}
\songpart{Intro}
\annotation{surf in the dark, a bottle cap on a table, felt piano}

\annotation{Couplet 1}\\
Deux bières sur la rambarde, la mer qui ne dort pas,\\
les feux d'un bateau lent, un avion qui s'en va.\\
Il m'a dit qu'il était né là-haut, à la montagne,\\
que la mer, depuis toujours, le fascinait tout bas.\\
Je l'ai regardé la regarder comme un enfant,\\
lui qui compte les sorties, lui qui se méfie du vent,\\
et pour une fois ses mains ne cherchaient rien,\\
posées sur le métal, tranquilles, sans plan.

\annotation{Refrain}\\
C'est dommage, c'est dommage,\\
il l'a dit doucement, sans me regarder.\\
Il m'a dit qu'il pourrait s'habituer à cette vie,\\
que c'était dommage que ce ne soit pas avec moi.\\
C'est dommage, c'est dommage,\\
et la mer a continué comme si de rien n'était.\\
J'ai souri, j'ai bu une gorgée,\\
j'ai fait semblant de ne pas comprendre pourquoi.

\annotation{Couplet 2}\\
Les lumières de Bal Harbour, sa deuxième maison,\\
ses chemises dans l'armoire, ses disques, ses raisons,\\
une vie bien rangée, à l'abri de la montagne,\\
où j'entre en invitée, le temps d'une chanson.\\
Il m'a dit, plus tard, que la montagne était loin,\\
que là-haut on n'entend pas la mer le matin.\\
J'ai pensé à une maison qui n'existera pas,\\
une fenêtre ouverte, et nous deux, et demain.

\annotation{Refrain}\\
C'est dommage, c'est dommage,\\
il l'a dit doucement, sans me regarder.\\
Il m'a dit qu'il pourrait s'habituer à cette vie,\\
que c'était dommage que ce ne soit pas avec moi.\\
C'est dommage, c'est dommage,\\
et la mer a continué comme si de rien n'était.\\
J'ai souri, j'ai bu une gorgée,\\
j'ai fait semblant de ne pas comprendre pourquoi.

\annotation{Pont}
\annotation{only bass, piano and the surf}\\
Je n'ai pas demandé avec qui, alors,\\
ni pourquoi pas, ni jusqu'à quand.\\
Il y a des phrases qu'on garde comme un coquillage,\\
on l'approche de l'oreille, et c'est la mer qu'on entend.\\
Un avion clignote et disparaît au nord,\\
la bière est tiède, et je l'aime encore.

\annotation{Refrain final}
\annotation{slower, accordion alone under the voice}\\
C'est dommage, c'est dommage,\\
il l'a dit doucement, et moi je l'ai cru.\\
Il pourrait s'habituer à cette vie,\\
moi aussi, moi aussi, mais je n'ai rien dit.\\
C'est dommage, c'est dommage,\\
la mer ne dort pas, et nous non plus.

\songpart{Outro}
\annotation{surf, a jet fading, piano holding one chord}\\
Né à la montagne…\\
la mer l'a toujours fasciné.\\
C'est dommage.
\annotation{End}
\end{lyrics}

